\documentclass[../../../include/open-logic-section]{subfiles}

\begin{document}

\olfileid{sth}{cardinals}{zfc}	
\olsection{$\ZFC$: یک نقطهٔ عطف}

با افزودن اصل خوش‌ترتیبی به آخرین نقطهٔ عطف نظری رسیده‌ایم. اکنون همهٔ
اصول موضوع لازم برای~$\ZFC$ را در اختیار داریم. به‌تفصیل:

\begin{defn}
نظریهٔ~$\ZFC$ دارای اصول موضوع زیر است: گسترش‌مندی، اجتماع، زوج‌ها،
مجموعه‌های توانی، نامتناهی، بنیاد، خوش‌ترتیبی و همهٔ نمونه‌های طرحواره‌های
جداسازی و جایگزینی. به‌بیانی دیگر، $\ZFC$ اصل خوش‌ترتیبی را به~$\ZF$
می‌افزاید.
\end{defn}

$\ZFC$ مخفف نظریهٔ مجموعه‌های \emph{زرملو--فرانکل} همراه با
\emph{انتخاب} است. شاید این امر اندکی عجیب به نظر برسد، زیرا اصلی که
افزودیم «خوش‌ترتیبی» نام داشت، نه «انتخاب». اما هنگامی که بعدتر اصل
انتخاب را صورت‌بندی کنیم، معلوم خواهد شد خوش‌ترتیبی (نسبت به~$\ZF$) با
انتخاب هم‌ارز است (بنگرید به~\olref[choice][woproblem]{thmwochoice}). پس
اینکه کدام‌یک را اصل «پایه» بگیریم تفاوتی ندارد. نام «$\ZFC$» نیز در
منابع کاملاً معیار است.

\end{document}
